\documentclass[12pt,a4paper]{article}
\usepackage[utf8]{inputenc}
\usepackage[T1]{fontenc}
\usepackage{placeins}
\usepackage{caption}
\usepackage[brazil]{babel}
\usepackage{amsmath} % para fórmulas matemáticas
\usepackage{amssymb} % para fórmulas matemáticas
\usepackage{geometry}
\usepackage{graphicx}
\usepackage{mdframed} % para spotbox
\geometry{
    a4paper,
    left=25mm,
    right=25mm,
    top=20mm,
    bottom=20mm
}
\newcommand{\customtitlepage}{
    \begin{titlepage}
        \vspace*{-0.4cm}
        \begin{figure}[!h]
            \centering
            \includegraphics[width=0.5\linewidth]{imagens/logoUFC.png}
        \end{figure}
        \centering
        {\large \textbf{Estatística para Engenharia}\par}
        \vspace{6cm}
        {\large \textbf{HOMEHORK}\par}
        {\large \textbf{Estatística e probabilidade}\par}
        \vspace{5cm}

        \begin{flushleft}
            \large Dupla: 
            \\ João Pedro Rolim Ximenes \\
            João Victor de Abreu Silva - 568274\\
            Profª: Michela Mulas \\
        \end{flushleft}
\vspace{3cm}
        \begin{flushleft}
            \centering \large \textbf{23/10/2025}
        \end{flushleft}
        \vfill
    \end{titlepage}
}

\begin{document}

\customtitlepage

\section{Questão 1}

 As emissões diárias de um gás poluente de uma planta industrial foram registradas 80 vezes, em uma determinada unidade de medida. Os dados obtidos estão apresentados na
 Tabela 1
 \begin{table}[!h]
\centering
\begin{tabular}{rrrrrrrrrrrr}
15.8 & 22.7 & 26.8 & 19.1 & 18.5 & 14.4 & 8.3 & 25.9 & 26.4 & 9.8 & 21.9 & 10.5 \\
17.3 & 6.2  & 18.0 & 22.9 & 24.6 & 19.4 & 12.3 & 15.9 & 20.1 & 17.0 & 22.3 & 27.5 \\
23.9 & 17.5 & 11.0 & 20.4 & 16.2 & 20.8 & 20.9 & 21.4 & 18.0 & 24.3 & 11.8 & 17.9 \\
18.7 & 12.8 & 15.5 & 19.2 & 13.9 & 28.6 & 19.4 & 21.6 & 13.5 & 24.6 & 20.0 & 24.1 \\
9.0  & 17.6 & 25.7 & 20.1 & 13.2 & 23.7 & 10.7 & 19.0 & 14.5 & 18.1 & 31.8 & 28.5 \\
22.7 & 15.2 & 23.0 & 29.6 & 11.2 & 14.7 & 20.5 & 26.6 & 13.3 & 18.1 & 24.8 & 26.1 \\
7.7  & 22.5 & 19.3 & 19.4 & 16.7 & 16.9 & 23.5 & 18.4 &      &      &      &      \\
\end{tabular}
\caption{Emissões diárias de gás poluente (questão 1).}
\end{table}

\subsection{ Calcule as medidas de tendência central (média, mediana e moda) e as medidas de
 dispersão (amplitude, variância, desvio padrão e coeficiente de variação) para o conjunto de dados da Tabela 1. Interprete os resultados} \\
Para facilitar a leitura da questão e dos dados em análise, é possível que seja criado um algoritmo em R utilizando o Sort que organize os números de maneira crescente, assim, fica mais fácil de interpretar os dados. Após essa implementação, a tabela obtida é: 

\begin{table}[!h]
\centering
\begin{tabular}{rrrrrrrrrr}
6.2 & 7.7 & 8.3 & 9.0 & 9.8 & 10.5 & 10.7 & 11.0 & 11.2 & 11.8 \\
12.3 & 12.8 & 13.2 & 13.3 & 13.5 & 13.9 & 14.4 & 14.5 & 14.7 & 15.2 \\
15.5 & 15.8 & 15.9 & 16.2 & 16.7 & 16.9 & 17.0 & 17.3 & 17.5 & 17.6 \\
17.9 & 18.0 & 18.0 & 18.1 & 18.1 & 18.4 & 18.5 & 18.7 & 19.0 & 19.1 \\
19.2 & 19.3 & 19.4 & 19.4 & 19.4 & 20.0 & 20.1 & 20.1 & 20.4 & 20.5 \\
20.8 & 20.9 & 21.4 & 21.6 & 21.9 & 22.3 & 22.5 & 22.7 & 22.7 & 22.9 \\
23.0 & 23.5 & 23.7 & 23.9 & 24.1 & 24.3 & 24.6 & 24.6 & 24.8 & 25.7 \\
25.9 & 26.1 & 26.4 & 26.6 & 26.8 & 27.5 & 28.5 & 28.6 & 29.6 & 31.8 \\
\end{tabular}
\caption{Emissões diárias de gás poluente em ordem crescente utilizando sort (questão 1).}
\end{table}
É possível responder o problema utilizando funções já conhecidas do R, como é o caso do mean e median, que já calculam automaticamente estas informações. No entanto, para melhor explanação, as fórmulas utilizadas são as seguintes:\\
Média (Soma de todos os elementos e divisão pela quantidade n deles):
\begin{equation} 
  \bar{x} = \frac{\sum_{i=1}^{n} x_i}{n} %para média aritmética
\end{equation}
Mediana (quando o resultado é um número quebrado, faz-se a média entre os dois valores):
\begin{equation}
    P = \frac{n + 1}{2}
  \end{equation}
Para a moda deve-se contar a quantidade de vezes que cada elemento aparece, e o que mais se repetir torna-se a moda daquele espaço amostral. \\
Amplitude (é a relação entre o valor máximo e mínimo de um determinado conjunto de resultados):
 \begin{equation}
    A = x_{\max} - x_{\min}
  \end{equation}
  Variância amostral (é feito a partir do somatório dos quadrados dos desvios simples sobre a quantidade elementos, neste caso):
  \begin{equation}
    s^2 = \frac{\sum_{i=1}^{n} (x_i - \bar{x})^2}{n}
  \end{equation}
  Desvio padrão (de forma sucinta, é a raíz da variância):
    \begin{equation}
    s = \sqrt{s^2} = \sqrt{\frac{\sum_{i=1}^{n} (x_i - \bar{x})^2}{n - 1}}
  \end{equation}
  Coeficiente de variação (se trata da relação percentual entre a variância e a média):
  \begin{equation}
    CV = \frac{s}{\bar{x}} \times 100\%
  \end{equation}
Com todos os cálculos devidamente apresentados, o resultado em R que calcula todos automaticamente foi a seguinte: \\
=== Medidas de Tendência Central ===\\
Média   : 19.02 \\
Mediana : 19.15 (lembrando que mediana é a média dos elementos entre as pos 40 e 41) \\
Moda    : 19.4\\ \\
=== Medidas de Dispersão ===\\
Amplitude             : 25.6 \\
Variância             : 30.84\\ 
Desvio padrão         : 5.55 \\
Coeficiente de variação: 29.2 porcento \\

\end{document}
